\begin{table}[t]
\fontsize{6.0pt}{7.2pt}\selectfont
\begin{tabular*}{\linewidth}{@{\extracolsep{\fill}}lcccc}
\toprule
Region & Fire Size km²,
mean (SD) & WUI
(\% in region) & Fatalities per 100K
(total) & Population Affected Annually, mean (SD) \\ 
\midrule\addlinespace[2.5pt]
Nationwide & 44.9 (189.6) & 3,652 (61.8\%) &    0.159 & 551,909 (1,049,594) \\ 
Alaska & 144.3 (384.0) &     1 (2.8\%) &    0.000 & 3,300 (4,234) \\ 
California & 94.6 (320.6) &   584 (71.6\%) &    0.577 & 2,182,606 (1,712,762) \\ 
Eastern & 2.8 (21.0) &   371 (57.4\%) &    0.034 & 386,120 (331,967) \\ 
Hawaii & 24.2 (38.8) &    15 (83.3\%) &    6.734 & 98,789 (103,394) \\ 
Intermountain & 109.4 (253.5) &   155 (47.0\%) &    0.240 & 207,921 (224,649) \\ 
Northern Rockies & 92.2 (187.2) &    91 (40.6\%) &    0.179 & 29,655 (31,541) \\ 
Pacific Northwest & 129.0 (251.4) &   217 (63.6\%) &    0.078 & 138,568 (202,655) \\ 
Southern & 12.4 (109.7) & 1,948 (64.6\%) &    0.066 & 1,668,933 (1,257,520) \\ 
Southern Rockies & 61.9 (128.0) &   145 (52.9\%) &    0.216 & 172,095 (328,366) \\ 
Southwestern & 119.4 (310.0) &   125 (59.0\%) &    0.227 & 107,885 (139,959) \\ 
\bottomrule
\end{tabular*}
\caption{Counts of wildfire burn zone disasters, structures destroyed, and civilian fatalities by modified US Forest Service region, 2000--2024. Percentages represent each region's share of the national total.}
\label{tab:fire_characteristics}
\end{table}

